\section{Introduction}
\label{sec:intro}

Resolved multi-wavelength surveys of nearby star-forming galaxies, most
prominently the Physics at High Angular resolution in Nearby GalaxieS
(\PHANGS) survey \citep{2021ApJS..255...19L,2023AJ....166..145L},
deliver aperture-level joint constraints on the cold gas, stellar, and
star-forming components of the interstellar medium (ISM) across
$\sim80$ galaxies and more than two orders of magnitude in environment.
These observations have established empirical scaling relations between
the local dynamical equilibrium pressure $\PDE$, the molecular gas
fraction $\fmol$, and the star formation rate surface density $\Ssfr$
\citep{2020ApJ...892..148S,2023ApJ...945L..19S,2025ApJ...989...66V},
which are reproduced in physical terms by the pressure-regulated,
feedback-modulated (\PRFM) framework \citep{2022ApJ...936..137O}.
\PRFM\ replaces a single empirical Kennicutt--Schmidt-type prescription
with a small set of physical parameters---feedback yields, effective
velocity dispersion, phase fractions---that connect resolved ISM
observables to the microphysics of feedback. The central inference
problem this paper addresses is the inverse one: given a vector of
\PHANGS\ aperture observables, what is the posterior distribution over
the physical parameters that govern ISM behavior, jointly with the
environmental parameters that describe the local galactic context?

This inverse problem is hard for a reason that is structural rather
than incidental. The forward map from environmental and physical
parameters to observed aperture summaries is mediated by a fully
nonlinear, multiphase, turbulent system whose snapshot realizations are
stochastic. Even at exact input parameters, the observable output of an
\TIGRESS-class simulation is a single draw from a distribution whose
moments---let alone whose full likelihood---are not analytically
accessible. The likelihood $p(\bm{x}_{\rm obs} \mid \bm{\theta})$
required for Bayesian inference therefore has no closed form, and
likelihood-based MCMC is not available. This is the canonical setting
of \emph{simulation-based inference} (\SBI), also termed
likelihood-free inference, in which the unknown likelihood is replaced
by an implicit one defined by a stochastic simulator
\citep{2020PNAS..11730055C}.

Modern \SBI\ trains a neural density estimator to approximate either
the posterior $p(\bm{\theta} \mid \bm{x}_{\rm obs})$ directly (neural
posterior estimation), the likelihood $p(\bm{x}_{\rm obs} \mid
\bm{\theta})$ (neural likelihood estimation), or the likelihood ratio
between competing parameter values (neural ratio estimation), using
$(\bm{\theta}, \bm{x}_{\rm sim})$ pairs drawn from the simulator
\citep{2020PNAS..11730055C}. Once trained, the estimator amortizes the
inference: posteriors can be evaluated at arbitrary observed data
without re-running the simulator. \SBI\ has been applied at scale to
cosmological parameter inference, galaxy-formation physics, and
gravitational-wave parameter estimation \CGK{(citations to add: Alsing+19,
Cranmer+20, Dax+21, Hahn+22 SimBIG, Villaescusa-Navarro+23 CAMELS-SBI;
verify ADS keys before inserting).} In each case, the strategy
succeeds when a structured simulation ensemble covers the parameter
volume densely enough to train the estimator, and when the summary
statistics chosen preserve the information that distinguishes
parameter values. Applications of \SBI\ to the resolved, multiphase,
star-forming ISM have so far been limited; this paper, to our
knowledge, is the first to deploy \SBI\ on \PHANGS\ aperture
observables using a structured \TIGRESS\ simulation suite as the
simulator. \CGK{Verify the novelty claim; soften to "among the first"
if a closely comparable effort exists.}

The simulation suite required for this program is the subject of the
companion paper \CGK{(\citealt{CITATION_NEEDED_TBD}; Paper~I, hereafter
\citetalias{CITATION_NEEDED_TBD}).} \citetalias{CITATION_NEEDED_TBD}
constructs a \PHANGS-informed simulation design for the \TIGRESS-NCR\
framework \citep{2024ApJ...972...67K}: a kernel density estimator is
fit to the joint \PHANGS\ aperture distribution of five environmental
parameters---the gas surface density $\Sgas$, stellar surface density
$\Sstar$, stellar scale height $\Hstar$, orbital frequency $\Omrot$,
and shear parameter $\qshear$---and a scrambled Sobol quasi-random
sequence is mapped through its marginal quantile functions to generate
the design. From each simulation we extract summary statistics that
correspond directly to \PHANGS\ aperture observables: $\Ssfr$, $\fmol$,
$\PDE$, the effective vertical velocity dispersion $\seff$, and the
thermal and non-thermal feedback yields $\Yth$ and $\Ynonth$. The
target of inference in this first paper is the posterior
$p(\thetaenv \mid \bm{x}_{\rm obs})$ over the five-dimensional
environmental parameter vector, conditioned on fiducial physics
($\Zg = 1$, $\Zd = 1$, standard IMF, calibrated \TIGRESS-NCR\
feedback). The framework is designed so that the physics parameter
vector $\thetaphys$ can be appended once the physics-extension block
described in \citetalias{CITATION_NEEDED_TBD} is run.

The scope of this paper is the construction, validation, and first
application of the inference framework. We present (i)~the choice of
summary statistics and the architecture of the neural posterior
estimator; (ii)~training on the environmental simulation suite of
\citetalias{CITATION_NEEDED_TBD}; (iii)~posterior recovery and coverage
diagnostics on held-out simulations and on mock observations drawn
from the simulator; and (iv)~application to \PHANGS\ aperture
measurements to infer $\thetaenv$ posteriors for each aperture and to
test the consistency of inferred parameters across galaxies and
across radial bins. Two extensions are deliberately deferred to future
work. The first is joint inference over $\thetaenv$ and $\thetaphys$,
which requires the physics-variation block of the suite. The second is
field-level inference using resolved \PHANGS\ maps, which requires the
synthetic observation pipeline outlined in
\citetalias{CITATION_NEEDED_TBD} together with non-Gaussian summary
statistics such as wavelet scattering coefficients
\citep{CITATION_NEEDED_TBD}. Both extensions are designed to slot into the
present framework without architectural changes.
\CGK{Confirm this is the intended scope for Paper~II; in particular,
whether posterior recovery on mock data and the \PHANGS\ application
are both in scope for this paper or whether the \PHANGS\ application
is itself deferred.}

The remainder of this paper is organized as follows.
\CGK{\autoref{sec:data}} describes the \PHANGS\ aperture observables
used as inference targets and the joining procedure between the
megatable and matched ancillary catalogs. \CGK{\autoref{sec:simulation_prior}}
summarizes the \TIGRESS-NCR\ simulation suite of
\citetalias{CITATION_NEEDED_TBD} as it is used here---the implicit prior
$p(\thetaenv)$, the simulator $p(\bm{x}_{\rm sim} \mid \thetaenv)$, and
the choice of summary statistics. \CGK{\autoref{sec:inference}}
presents the inference framework: the neural posterior estimator, the
training procedure, and the coverage and calibration diagnostics used
to validate it. \CGK{\autoref{sec:results}} presents the inferred
posteriors on held-out simulations and on \PHANGS\ apertures, and
discusses the implications for ISM physical parameters. We conclude
in \CGK{\autoref{sec:conclusions}.}
\CGK{Section labels assumed from the existing sections/ directory
listing (data, simulation\_prior, inference, results, conclusions);
adjust if the final structure differs.}
